\section{Homogeneous cuts containing a prescribed central closed subscheme}
Let $k$ be an infinite field, let $R$ be a commutative $k$-algebra,
and put $B=R[X_j\mid j\in\sigma]$ with its actual polynomial grading.
The specialization $R=k[t]$, with finitely many original projective
coordinates, retains the actual polynomial ambient space of a projective
family. The statements below construct original homogeneous equations;
no surface, curve family, ray or norm comparison is an input.

\begin{lemma}[A homogeneous ideal has a homogeneous nonvanishing witness]
\label{lem:pl-gauss-homogeneous-ideal-nonvanishing-witness}
\lean{EnvelopeProperty.PLGaussVerticalCurveSurfaceSelection.exists_homogeneous_mem_notMem}
\leanok
Let $I$ be an actual homogeneous ideal of $B$ and $P$ an ideal with
$I\not\subseteq P$. There are $e\ge0$ and an original homogeneous
$g\in I_e$ such that $g\notin P$.
\end{lemma}
\begin{proof}
Choose $a\in I\setminus P$ and write its finite sum of actual homogeneous
components $a=\sum_e a_e$. Homogeneity of $I$ puts each $a_e$ in $I$.
If every component belonged to $P$, closure of $P$ under finite sums
would put $a$ in $P$. Some $a_e$ therefore does not belong to $P$;
this component is the required $g$ of its actual original degree.
\end{proof}

\begin{theorem}[Large homogeneous equations avoid finitely many relevant primes]
\label{thm:pl-gauss-eventual-homogeneous-prime-avoidance}
\uses{lem:pl-gauss-homogeneous-ideal-nonvanishing-witness}
\lean{EnvelopeProperty.PLGaussVerticalCurveSurfaceSelection.exists_eventually_homogeneous_avoiding}
\leanok
Let $(P_i)_{i\in T}$ be a finite family of actual prime ideals of $B$.
Assume $I\not\subseteq P_i$ and that for each $i$ some actual coordinate
$X_{j_i}$ does not belong to $P_i$. There is $N\ge0$ such that for EVERY
$d\ge N$ there is an original homogeneous $f\in I_d$ with
$f\notin P_i$ for every $i\in T$. The empty family is included.
\end{theorem}
\begin{proof}
For each $i$, choose $g_i\in I_{e_i}\setminus P_i$ by the preceding
lemma, and choose $X_{j_i}\notin P_i$. Let $N$ be the maximum of the
finitely many $e_i$, with maximum zero for the empty family. For $d\ge N$,
the actual polynomial
\[
 h_i=g_iX_{j_i}^{d-e_i}
\]
has degree $d$, belongs to $I$, and does not belong to $P_i$: primality
and avoidance of both factors exclude membership, including when
$d=e_i$. Consider the actual $k$-vector space $W_d=I\cap B_d$.
The set $W_d\cap P_i$ is a $k$-linear subspace and is proper because it
does not contain $h_i$. A finite union of proper linear subspaces of a
vector space over an infinite field is proper. Apply this fact to the
actual subspaces $W_d\cap P_i$ and choose $f$ outside their union.
It belongs to the same original ideal and homogeneous degree space,
and avoids every original prime. For the empty family the zero element
of $W_d$ satisfies the stated conclusion.
\end{proof}

\begin{lemma}[Every original polynomial projective point has a nonzero coordinate]
\label{lem:pl-gauss-original-projective-point-coordinate}
\lean{EnvelopeProperty.PLGaussVerticalCurveSurfaceSelection.exists_coordinate_notMem}
\leanok
Every actual point $x\in\operatorname{Proj}B$ has an actual coordinate
$X_j\notin\mathfrak p_x$.
\end{lemma}
\begin{proof}
Suppose every coordinate belonged to $\mathfrak p_x$. Every positive
degree homogeneous polynomial is a sum of monomials of positive degree.
Each monomial has a coordinate factor, so every such polynomial belongs
to the ideal generated by the coordinates and hence to $\mathfrak p_x$.
The irrelevant ideal is generated by these positive degree pieces.
It would therefore be contained in $\mathfrak p_x$, contradicting the
actual relevance condition defining the original projective point.
\end{proof}

\begin{corollary}[Actual projective hypersurfaces contain the prescribed locus]
\label{cor:pl-gauss-original-projective-hypersurface-avoidance}
\uses{thm:pl-gauss-eventual-homogeneous-prime-avoidance,lem:pl-gauss-original-projective-point-coordinate}
\lean{EnvelopeProperty.PLGaussVerticalCurveSurfaceSelection.exists_eventually_projective_hypersurface_avoiding}
\leanok
Let $x_i$ be a finite family of actual points of $\operatorname{Proj}B$
outside the original closed locus $V_+(I)$. There is $N\ge0$ such that
for every $d\ge N$ an actual original homogeneous $f\in I_d$ defines a
hypersurface $V_+(f)$ containing $V_+(I)$ and avoiding every $x_i$.
This applies to the original polynomial projective ambient space over
$k[t]$, with no supplied section or hypersurface certificate.
\end{corollary}
\begin{proof}
The actual projective points supply their own prime homogeneous ideals.
Being outside $V_+(I)$ says exactly that $I$ is not contained in each
of these actual ideals. The preceding coordinate lemma internally
supplies the relevance witnesses. Apply the homogeneous avoidance
theorem to obtain $f$. Since $f\in I$, every prime containing $I$
contains $f$, so $V_+(I)\subseteq V_+(f)$. The inequalities
$f\notin\mathfrak p_{x_i}$ say exactly that each original $x_i$ lies
outside this actual hypersurface. All points and coefficients are the
unchanged original ones.
\end{proof}

These equations provide the first cut used to retain a prescribed
central curve while avoiding the generic points of larger central
components. They do not by themselves prove the existence of a
horizontal surface, the exact dimension of a cut component, preservation
of the prescribed curve through repeated cuts, or its finite dominant
lift after normalization. Those geometric conclusions require their
own actual dimension, flatness and normalization arguments.
